\documentclass{rvcepaper}
\begin{document}

% ================= BLOCK A: 2025-26 papers (inserted after the Apr-Jun 2026 block)
% ---------- CV266TEE Intelligent Transport Systems, CIE - I, 16-04-2026
\begin{xpaper}
\paperinfo{deptline={Department of Civil Engineering},date={16-04-2026},exam={CIE -- I},maxmarks={60},sem={VI},prog={UG},duration={2 Hrs},
 title={Intelligent Transport Systems},code={CV266TEE}}
\xhead
\ptitle{Quiz}
\begin{xtable}{Questions}
\q{1.}{List any two functions of emergency management.}{2}{2}{1}
\q{2.}{List four stages involved in Transportation Planning process.}{2}{1}{1}
\q{3.}{State the goal of Intelligent Transportation Systems (ITS).}{2}{1}{1}
\q{4.}{List any two aims of personalized public transit.}{2}{2}{2}
\q{5.}{List any two functions of maintenance management.}{2}{2}{2}
\end{xtable}
\ptitle{Test}
\begin{xtable}{Questions}
\q{1.}{Explain the cycle of urban transport that leads to congestion on the road network.}{10}{3}{1}
\q{2.}{Explain what is meant by an ITS user service and its basic characteristics.}{10}{3}{2}
\q{3.}{Explain why training in Intelligent Transportation Systems (ITS) is essential and identify educational requirements.}{10}{3}{1}
\q{4.}{Explain in detail about Travel and traffic management user service.}{10}{3}{1}
\q{5.}{Identify the busiest corridor in your vicinity. What challenges does this corridor face, and how can technology address these issues?}{10}{4}{2}
\end{xtable}
\end{xpaper}

% ---------- MA266TEU Mathematical Modelling, CIE - 1, 16/04/2026
\begin{xpaper}
\paperinfo{deptline={Department of Mathematics},date={16/04/2026},exam={CIE - 1},maxmarks={100},sem={VI},prog={UG},duration={2 Hrs. (2.30pm-4.30pm)},
 title={Mathematical Modelling},code={MA266TEU}}
\xhead
\begin{xtable}{Quiz: Questions}
\q{1}{The maximum horizontal distance covered by a body of mass 5 kg is projected from the ground at an angle of $45^\circ$ with an initial velocity of $20\,m/s$ is \blank[1.2cm]. (Take acceleration due to gravity as $g \approx 9.8\,m/s^2$)}{2}{L1}{1}
\q{2}{The growth in the number of software engineers in a country follows the law $\dfrac{dx}{dt} = \dfrac{1}{20}x$. If initially there were 800 engineers, then after 20 years, their approximate number will be \blank[1.2cm] and their doubling period is approximately \blank[1.2cm] years.}{2}{L1}{2}
\q{3}{A heated metal sphere at $120^\circ$C is placed in a room maintained at $30^\circ$C. If the cooling constant is $k = 0.03$, then the temperature of the sphere after 15 minutes is \blank[1.2cm].}{2}{L2}{3}
\q{4}{The sequence 2, 5, 10, 17, 26, \dots can be modelled by a difference equation of the form \blank[1.2cm].}{2}{L2}{2}
\q{5}{State Fick's law of diffusion model.}{2}{L1}{1}
\end{xtable}
\ptitle{Test}
\begin{longtable}{|C{1.05cm}|p{12.2cm}|C{0.95cm}|C{0.95cm}|C{0.95cm}|}
\hline \textbf{SL. No} & \multicolumn{1}{c|}{\textbf{Test: Questions}} & \textbf{M} & \textbf{BT} & \textbf{CO}\\\hline
\q{1a.}{Model the motion of a projectile launched in a vacuum with an initial velocity at a given horizontal angle. Formulate the governing equations of motion, solve them and determine key characteristics such as time of flight, maximum height and horizontal range.}{6}{L3}{3}
\q{1b.}{Discuss how mathematical modelling is classified.}{4}{L1}{2}
\q{2.}{A nutrient is periodically injected into a patient's bloodstream at fixed intervals of time $T$. Meanwhile, the nutrient is continuously absorbed by the body at a rate proportional to its concentration $c(t)$ in the blood at time $t$. Derive the differential equation governing $c(t)$. Solve it and interpret the long-term behaviour of the nutrient concentration in the bloodstream.}{10}{L4}{4}
\q{3.}{Develop a differential equation model based on the logistic law of population growth, solve it and analyse any one real life scenario.}{10}{L3}{2}
\q{4.}{Describe the systematic 12-step procedure for tackling a problem through the use of mathematical modelling.}{10}{L1}{1}
\q{5a.}{Solve the difference equation $x_{n+1} - 0.3x_n - 6 = 0$, $x_0 = 10$, $n \ge 0$. Hence find the limiting (convergent) value of the sequence.}{5}{L3}{3}
\q{5b.}{Form the difference equation for the sequence 0, 1, 1, 2, 3, \dots and solve it.}{5}{L2}{2}
\end{longtable}
\end{xpaper}

% ---------- HS251TA Principles of Management and Economics, CIE - I, 06.11.2025
\begin{iempaper}
\paperinfo{shownote=0,date={06.11.2025},exam={CIE -- I},maxmarks={10 + 50},sem={V},prog={UG},duration={30 + 90 Min},
 title={Principles of Management and Economics},code={HS251TA}}
\iemhead
{\small Note:}\par\vspace{-1pt}
\hspace*{0.5cm}{\footnotesize 1. \hspace{0.1cm}Answer all the Questions.}\par\vspace{3pt}
\begin{cietable}{M}{BT}{CO}
\qpart{Part -- A}
\q{1}{Why do managers need to be both efficient and effective?}{2}{L1}{01}
\q{2}{Mention three important skills required to be an effective manager.}{2}{L2}{02}
\q{3}{Differentiate between Unity of command and Unity of direction.}{2}{L2}{01}
\q{4}{Why should managers have a basic knowledge of economics?}{2}{L2}{04}
\q{5}{What does the circular flow of income illustrate in an economy?}{2}{L1}{04}
\qpart{Part -- B}
\q{1}{A manufacturing firm is facing issues with productivity and frequent machine breakdowns. The management decides to reorganize shop-floor activities and assign specific tasks to workers based on time and motion studies. Incentives are planned for those who achieve higher productivity. Analyse the above case using Frederick Taylor's Scientific Management principles. Which principles are being applied, and how could they help improve efficiency?}{10}{L2}{01}
\q{2}{A growing IT company notices that despite excellent pay, employees show low morale and poor teamwork. The management brings in a consultant who focuses on employee well-being, informal group interactions, and feedback sessions. Evaluate the management's approach using Elton Mayo's Human Relations Theory. What lessons can be learned from the Hawthorne Experiments in this context?}{10}{L3}{02}
\q{3}{A CEO represents the organization at a global conference, negotiates mergers, sets strategic goals, and ensures communication across all divisions. Explain how this situation reflects Mintzberg's 10 managerial roles, grouped into their respective categories.}{10}{L3}{01}
\q{4}{A leading automobile manufacturer notices that the demand for its vehicles has declined due to an increase in petrol prices and the entry of a new competitor offering electric vehicles at lower prices. Meanwhile, the government is implementing policies to stabilize inflation and boost employment. Using the principles of microeconomics and macroeconomics, analyse how both these levels of economic study help in understanding and addressing this situation.}{10}{L2}{04}
\q{5}{Explain the two-sector circular flow model with the help of an example. How does it demonstrate the interdependence between households and firms in an economy?}{10}{L2}{04}
\end{cietable}
\end{iempaper}

% ================= BLOCK B: 2024-25 odd-semester paper (appended at the end of the file)
% ---------- HS251TA Principles of Management and Economics, CIE - I, 25th November 2024
\begin{iempaperB}
\paperinfo{ayear={2024-2025 (ODD Sem)},date={25\textsuperscript{th} November 2024},code={HS251TA},sem={V},exam={CIE -- I},maxmarks={10 + 50},duration={120 Min},
 title={PRINCIPLES OF MANAGEMENT AND ECONOMICS}}
\iemheadB
\begin{cietable}{M}{BT}{CO}
\qpart{Part -- A}
\q{1}{Which step in POSDCORB involves delegating tasks and giving structured instructions?}{01}{1}{1}
\q{2}{Define ``effectiveness'' in the context of management functions.}{01}{1}{1}
\q{3}{Which management approach is most associated with using data and mathematical models to improve decision-making?}{01}{1}{1}
\q{4}{In the Contingency Approach, what determines the best management style?}{01}{1}{2}
\q{5}{\blank[2cm] is the primary scientist associated with the foundation of Administrative Theory.}{01}{1}{1}
\q{6}{How Total Quality Management (TQM) incorporates the Quantitative approach to improve customer satisfaction.}{01}{1}{2}
\q{7}{What role does fiscal policy play in stabilizing an economy during a recession?}{01}{1}{4}
\q{8}{What is the role of Government within the circular flow model?}{01}{1}{4}
\q{9}{What are the key components of macroeconomics that measure the health of an economy?}{02}{1}{5}
\qpart{Part -- B}
\q{1a}{Mention the underlying reasons for the increase in productivity during the Hawthorne Studies, and how do these reasons challenge the idea that only physical conditions affect workplace efficiency?}{06}{4}{1}
\q{1b}{How does Contingency Theory explain the limitations of a `one-size-fits-all' approach to management? Provide two examples to illustrate how it applies to different organizational situations.}{04}{3}{2}
\q{2}{Compare and contrast the Classical and Behavioral approaches to management. How do these approaches address productivity, and what are their limitations in modern workplaces?}{10}{2}{1}
\q{3}{Analyze how POSDCORB ensures both efficiency and effectiveness in management. Discuss the potential challenges managers face if one of these functions is underperformed.}{10}{3}{1}
\q{4}{Evaluate the relevance of Systems Theory in managing modern organizations. How does its emphasis on interdependence and the environment benefit decision-making?}{10}{4}{1}
\q{5}{Analyze how microeconomics, macroeconomics, and economic systems interact within the framework of the circular flow model. What challenges arise when these components fail to work in harmony?}{10}{4}{5}
\end{cietable}
\btfoot{|l|l|l|*{11}{C{0.8cm}|}}{\hline
\multirow{2}{*}{\shortstack{Marks\\Distribution}} & \multicolumn{2}{c|}{Particulars} & CO1 & CO2 & CO3 & CO4 & CO5 & L1 & L2 & L3 & L4 & L5 & L6\\\cline{2-14}
 & Max & Quiz & 04 & 02 & -- & 02 & 02 & 10 & -- & -- & -- & -- & --\\\cline{3-14}
 & Marks & Test & 36 & 04 & -- & 0 & 10 & -- & 10 & 14 & 26 & -- & --\\\hline}
\end{iempaperB}

\end{document}
